\section{Gameplay Basics}

In this section we will go over the basics on how to play and get a gist of the fantasy this book is trying to provide. It also works as glossery to where to find what you are looking for
\subsection{Rolls}

Any conflict, be it against a fellow cultivator, a spirit beast trying to take your families lives or the heavens themselves striking you down for defying fate itself. It is resolved through a dice roll. (Though laughably easy tasks or impossible (for now) tasks the GM may forgo a roll at all)

A basic dice roll consists of a base roll of 2d6 + a dice pool value from your cultivation + a skill value
The dice pool value is the highest number shown on a pool of d6. The number of dice is the difference between your level of cultivation and the obstacle. If the obstacle is of higher cultivation the number rolled will be subtracted instead. This can be negated by one step by doubling the number of cultivators doing the task
The skill value is the difference between your skill at the task and the skill/difficulty of your obstacle

\begin{itemize}
    \item If the total is below 6 you fail
    \item If the total is between 6 and 10 you can succeed at a great cost
    \item If the total is between 10 and 12 you can succeed at a minor cost
    \item If the total is above 12 you succeed without cost
\end{itemize}


(To come out of a task without having sacrificed anything you must have some sort of advantage)

\subsection{Combat}

When entering into combat every combatant rolls a speed roll against the lowest cultivation
combatants. Highest speed skill rank breaks ties, in case of equal skill ranks and rolls the player goes first.

Every combatant is then listed from first to last, everyone in order gets to move and do an action during their turn.
Once everyone has done their turns the first combatant gets to take their turn again, starting a new round.

If every combatant is on one side of a battle the measures can be normalised.

The location and distance between combatants are abstracted into measures. 
Far, Ranged, Long, Medium \& Short on each side of the battle.

A character can move between neigbouring measures once per turn. \\
One measure can be safely moved through if you forgo attacking.\\ % Step
One measure can be rashly moved through if you risk getting attacked.\\ % Stride 
Two measured can be rashly moved through if you both forgo attacking and risk getting attacked. % Dash

If a combatant moved rashly within, into or out of your weapons measure you can once per round attack that combatant.


An attack consists of hitting and dealing damage. 
First you roll an appropriate weapon roll against a combatants defence skill. For melee weapons its your weapon skill targeting the other combatants Parry skill.
If your weapon is ranged you roll against their Evade instead.

If you succeed, you then roll an appropriate weapon roll again but against the enemies Hardiness skill. 

If you succeed, you then deal a wound to them. 

If you succeed with a minor cost you may give your opponent a point of advantage or suffer a minor narativly applicable consequence.

If you succeed with a great cost you may give your opponent two points of advantage or suffer a Major narativly applicable consequence.

Advantage points can be spent to increase your relative skill against a combatant during an attack.

If you dealt a wound to a combatant you can deal another at the cost of two Advantage points.

\subsection{Statistics}

The most important aspect of a cultivator is their Cultivation Stage and their Qi

Your Cultivation Stage determines your position, respect and access in the cultivation world. It determines capacity to collecting of ambient Qi which furthers your progress. 
As a rule of thumb the required effort to move between stages doubles every two stages.

Your Qi is divided into your maximum Qi which is the amount your Dantian is currently able to hold without problems. Your Current Qi is how much you have available and circulating your body. The latter one is easier to replenish than to increase your maximum

Your Qi can also be elementally aligned to any of the five elements. If your elemental alignment is advantageous towards your task you can grow your own dice pool (d6->d8) or if it is disadvantageous you decrease your dice pool (1d6->1d4)

The next important one is your capacity for Wounds. You can harbor a number of Wounds equal to your Cultivation Stage * 2 + 1. Along with wounds usually comes ailments which there is no limit to how many you can gain but gaining duplicates will transfer into Wounds

Physical damage is not the only type a cultivator has to face in their journeys. Qi Deviation is the chaos within your Qi and is tracked by a percentage value. This is how much of your maximum stored Qi which is inaccessible. Qi Deviation can be dealt with by brute force by flushing it out with great amounts of Qi

More about this can be (Will be) read in the Cultivation section

\subsubsection{Defence}

To avoid taking damage a myriad of defences is needed


\textbf{Hardiness}: Your ability to shrug off hits which made contact\\
This skill comes into play when an enemy is trying to see whether a hitting attack turns into a wound. It is also apllicable as a roll when fighting off poisons and infections.

The base value of this skill starts at 3 and every rank taken is added ontop of this.

\textbf{Parry}: Your ability to redirect your enemies close quarter strikes\\
This skill comes into play when an enemy makes a melee attack of any kind against you. 

The base value of this skill starts at 3 and every rank taken is added ontop of this.


\textbf{Evade}: Your ability to evade your enemies long rage volleys\\
This skill comes into play when an enemy makes a ranged attack of any kind against you. 

The base value of this skill starts at 3 and every rank taken is added ontop of this.


\textbf{Stealth}: Your ability to pass unnoticed or disguised\\
This skill comes into play when you try to move silently or quickly slip past someone. 

If an advercary is unaware of you their detect roll must succeed without without cost to notice you.

If they are aware of of general threats but not you in perticular (an unalerted guard) their detect roll must succeed 
with atleast a minor cost in which case will lend you a single turn before they manage to react.

If they are searching for you in perticular their detect roll can succeed 
with a great cost in which case will lend you a single turn before they manage to react.

The base value of this skill starts at 6 and every rank taken is added ontop of this.

\textbf{Wits}: Your ability to see through deceptions, shrug off empathy \& ward off certain magics\\
This skill comes into play when someone tries to appeal to your emotions or fool you. It is also your defence against certain
types of magical effects on the mind.

An active Wits roll is rare since most of the Mental skills are more active.

The base value of this skill starts at 6 and every rank taken is added ontop of this.

\textbf{Resolve}: Your ability to stand against intimidation, temptation and others conviction\\
This skill comes into play when someone tries to intimidate, tempt or convince you. 

This skill can be used for rolls against succumbing to tempting flaws like Hedonism.

\subsubsection{Offence}

To deal with those pesky rival cultivators 


\textbf{Short weapons}: Your ability to use weapons with a measure of no more than a foot\\
This skill comes into play whenever you're trying to use any weapon with a short measure. 

If an enemy uses a weapon with a long measure your rank in this skill is increased by one.\\
If an enemy uses a weapon with medium measure your rank in this skill is decreased by one. 

Some of these weapons can also be thrown from up to a Long measure and will target the targets Evade instead of Parry.

\textbf{Medium weapons}: Your ability to use weapons with a measure of up to an arms\\
This skill comes into play whenever you're trying to use any weapon with a medium measure. 

If an enemy uses a weapon with a short measure your rank in this skill is increased by one.\\
If an enemy uses a weapon with long measure your rank in this skill is decreased by one. 

\textbf{length Long weapons}: Your ability to use weapons with a measure of up to your\\
This skill comes into play whenever you're trying to use any weapon with a long measure.

If an enemy uses a weapon with a medium measure your rank in this skill is increased by one.\\
If an enemy uses a weapon with short measure your rank in this skill is decreased by one. 

\textbf{Ranged weapons}: Your ability to use projectile based weapons. \\
This skill comes into play whenever you're trying to use any weapon with a ranged measure.

If an enemy is at medium measure of you your rank in this skill is decreased by two.\\
If an enemy is at short measure you can not make a ranged weapons roll against them.

\textbf{Unorthodox weapons}: Your ability to use weapons outside of any previous category. Skills does not transfer between different Unorthodox weapons
This skill comes into play whenever you're trying to use an unorthodox weapon. 

Specific range, modifiers \& rules are stated under each unorthodox weapon.

\subsubsection{Specialist}

Not every Cultivator is a crazed sword lunatic and at least dabble in some sort of specialisation

\textbf{Medicine}: Your ability to treat poisons, mend wounds, set broken bones into place\\
\textbf{Divination}: Your ability to read the stars, tea leaves or entrails to predict the future\\
\textbf{Meditation}: Your ability to calm your mind to treat spiritual maladies or stave off possession \\
\textbf{Talent}: Your ability to do anything not covered by any other specialist skill\\
\textbf{Trade}: Your ability to partake in the art of creation be it Alchemy, Smithing or Weaving\\
\textbf{Survival}: Your ability to sustain and navigate areas, ranging from Cities to deserts to caves\\
\textbf{Ritual}: Your ability to hold and perform special rituals with varying effects 

\subsubsection{Physical}

The world is perilous and roads are not all paved pavilions, a cultivator needs to be able to traverse it

\textbf{Athletics}: Your ability to perform physical tasks not covered by any other skill\\
\textbf{Swim}: Your ability to move and perform action while in more perilous waters\\
\textbf{Speed}: Your ability to move on ground but also react to the world around you.\\
\textbf{Muscle}: Your ability to lift, throw and break the world around you \\
\textbf{Endurance}: Your ability to endure harsh climates and push past your physical limits \\
\textbf{Ride}: Your ability to manoeuvre, calm and order a beast. Usually horses \\
\textbf{Sail}: Your ability to manage, steer and upkeep a water vessel.

\subsubsection{Knowledge}

Mysteries are abound in the world, but ones best weapon against it is insight on how it works.

\textbf{History}: Your ability to recall a region's past and its reasons \& consequences\\
\textbf{Creatures}: Your ability to recall or infer a type of creatures biology, weaknesses \& strengths\\
\textbf{Places/Cultures}: Your ability to understand customs, politics, food, traditions, etc. \\
\textbf{Martial Disc.}: Your ability to identify and categorise different martial arts and its nature\\
\textbf{Institutions}: Your ability to understand the inner workings of things like armies or sects.\\
\textbf{Language}: Your ability to speak and understand a language and the ones close by\\
\textbf{Read Script}: Your ability to read and write a language or sects internal script.\\
\textbf{Religion/Gods}: Your ability to participate, identify and differentiate different religions\\
\textbf{Classics}: Your ability to read, recite and comprehend classic literature.

\subsection{Allocation}

A character has two starting categories which they will start out better at than the other. 
At character creation a 12 points can be allocated into those catagories and 6 points in the rest
Whenever a character goes through a breakthrough in their cultivation they gain points to either spend or save. 

The cost to increase the rank of a skill increases in pase with its rank. 

\begin{tabular}{|c|c|c|l|}
    \hline
    Rank & Cost & Total* & Description\\
    \hline
    0 & 0 & 0  & \thead{Untrained}  \\
    1 & 1 & 1  & \thead{Novice} \\
    2 & 2 & 3  & \thead{Student} \\
    3 & 3 & 6  & \thead{Scholar} \\
    4 & 4 & 10 & \thead{Master} \\
    5 & 5 & 15 & \thead{Earthly \\Authority} \\
    6 & 6 & 21 & \thead{Celestial \\Authority} \\
    \hline
\end{tabular}

While being below the Nascent Soul stage you are unable to advance to Rank 4 in any one skill. 
The same Applies to XXX stage and Rank 5 and XXX stage and Rank 6 